\section{Introdu\c{c}\~ao}
\label{sec:introducao}

As audiências públicas das comissões da Câmara dos Deputados reúnem parlamentares, representantes
do governo, especialistas e entidades da sociedade civil em torno de propostas legislativas e
políticas públicas, e suas transcrições registram o que cada participante defendeu. Essas
transcrições têm em média 18 mil palavras no PublicHearingBR~\citep{fernandes2024publichearingbr}, e o
público costuma conhecê-las por matérias jornalísticas que têm em média 4,2\% do tamanho da
transcrição e citam poucas opiniões de poucos participantes. O conjunto de dados registra o que a
matéria atribui a cada pessoa, mas não o ponto da fala de onde cada opinião foi extraída.

Este trabalho trata de três perguntas encadeadas que partem dessa lacuna: qual trecho da fala de uma
pessoa sustenta cada opinião que a matéria lhe atribui; como representar, somente a partir da fala,
o que uma pessoa defende ao longo de várias audiências, de modo que a representação oriente um
modelo de linguagem; e se um LLM que recebe essa representação consegue estimar como a pessoa se
posicionaria numa audiência que não entrou nela, declarando a base da estimativa. Uma estimativa
desse tipo pode apoiar a preparação de audiências e ferramentas de transparência que informem o que
um participante já defendeu sobre um tema e como suas posições se aplicam a um tema novo. LLMs usados
para simular pessoas, porém, tendem a produzir respostas estereotipadas e a apagar a diversidade
interna dos grupos~\citep{bisbee2024synthetic,wang2025large}, e o modelo pode responder com a própria
opinião ou com o que associa ao nome da pessoa. Por isso, a simulação proposta usa somente as falas
da própria pessoa, exclui o partido do prompt, declara a base de cada estimativa e é avaliada com
opiniões publicadas sobre audiências posteriores às usadas no perfil.

A abordagem (Seção~\ref{sec:metodo}) tem três componentes: a Unidade Deliberativa Verificável (UDV),
que liga cada opinião do resumo estruturado a uma sentença da fala da própria pessoa, com offsets de
caracteres, um nível que declara o tipo de vínculo e a origem do vínculo; os perfis por ator, escritos
por um LLM somente com as falas de cada pessoa nas audiências de treino, em blocos de posições,
critérios e valores, alinhamentos declarados e forma de argumentar; e a simulação, que fornece o
perfil a um LLM, acrescenta trechos das falas anteriores da pessoa recuperados pelo assunto e usa
\textit{classifier-free guidance} para amplificar o efeito desse material. As UDVs das audiências de
validação e de teste servem de gabarito para perguntas de múltipla escolha que medem o efeito de cada
componente. O trabalho se inspira nas Trilhas A e B do desafio: a UDV e a recuperação restrita à fala
da pessoa tratam da auditoria de documentos longos, e os perfis e a simulação tratam da mineração de
opiniões e de atores, sempre com apoio em trechos localizados da transcrição.

\medskip\noindent\textbf{Contribui\c{c}\~oes.} As principais contribui\c{c}\~oes deste trabalho s\~ao:
\begin{itemize}
  \itemsep2pt
  \item a UDV, um registro que liga 2.105 das 2.203 opiniões do PublicHearingBR a uma sentença da
  fala da própria pessoa, com offsets, nível e proveniência, conferido por um verificador separado;
  \item perfis por ator escritos somente com falas de audiências de treino, com regras que impedem
  atribuir à pessoa o que vem do cabeçalho da audiência, de terceiros, do partido ou de conhecimento
  externo às falas;
  \item um método de simulação de participantes que combina o perfil, trechos recuperados da fala da
  própria pessoa e \textit{classifier-free guidance}, com classificação da base da estimativa e
  justificação conferida por código;
  \item um protocolo de avaliação por múltipla escolha com gabarito derivado das UDVs, distratores da
  mesma audiência, leitura pelas probabilidades das letras em quatro rotações das opções, escolha de
  hiperparâmetros na validação e intervalos por \textit{bootstrap} de audiências, aplicado às
  condições da simulação com <resultado principal>.
\end{itemize}
